\chapter{Multi-Class Classification Results}
\label{ch:multiclass}

\section{The Generalisation Collapse}

Table~\ref{tab:multiclass-headline} contains the central empirical claim
of the thesis. On the random stratified split, XGBoost achieves macro
F1 = 0.86 and random forest follows at 0.84. On the day-based split,
every tree ensemble collapses to a macro F1 in the 0.43--0.46 range. The
collapse is uniform across all three tree models, which differ by less
than 0.04 macro F1 on the day split. The reason: Friday's DDoS, PortScan,
and Bot classes never appear in Mon--Wed training, so the classifier
cannot output the correct label for those flows no matter how well it
is calibrated.

\begin{table}[H]
    \centering
    \caption{CICIDS2017 multi-class macro F1 by split policy.}
    \label{tab:multiclass-headline}
    \small
    \begin{tabular}{lccc}
        \toprule
        \textbf{Model} & \textbf{Stratified} & \textbf{Day-based} & \textbf{$\Delta$} \\
        \midrule
        LogReg & 0.266 & 0.429 & +0.163 \\
        Random Forest & 0.845 & 0.438 & $-$0.407 \\
        XGBoost & 0.863 & 0.440 & $-$0.423 \\
        LightGBM & 0.300 & 0.463 & +0.163 \\
        \bottomrule
    \end{tabular}
\end{table}

\interpretation{The 0.4-point drop for tree ensembles between split
policies is the single largest signal in this thesis. The published
numbers reflect in-distribution memorisation, not deployment readiness.}

\begin{figure}[H]
    \centering
    \includegraphics[width=0.95\textwidth]{generalisation_gap_all_models.png}
    \caption{Stratified versus day-split macro F1 across all four models.
    Blue bars are stratified (published-style); red bars are day-based
    (deployment-style). Every tree ensemble loses around 0.4 macro F1.}
    \label{fig:gen-gap}
\end{figure}

Logistic regression behaves differently: its macro F1 actually
\textit{improves} from 0.27 stratified to 0.43 on the day split. Its
stratified score was already low because it cannot capture the non-linear
class boundaries that tree models exploit. On the day split, the unseen
classes hurt all models equally, but the tree models also lose the
spurious correlations they had memorised, while logistic regression had
not learned those correlations to begin with.

\section{Why Multi-Class Collapses: Confusion Analysis}

The day-split confusion matrix shows the failure mode. Friday DDoS flows
are predominantly predicted as DoS Hulk, the closest seen class in
feature space. PortScan splits between DoS Hulk and Benign. Bot is
mostly predicted as Benign, because without seeing C2 traffic in
training the model cannot distinguish a periodic beacon from background
process traffic. Every Friday error is structural --- a consequence of a
label not existing in training rather than a tunable model failure ---
and it cannot be addressed by additional training on the same days.

\begin{figure}[H]
    \centering
    \includegraphics[width=0.95\textwidth]{confusion_matrix.png}
    \caption{Confusion matrices: stratified split (left) versus day split
    (right) for the best multi-class model. On stratified the diagonal
    is bright everywhere; on day the Friday rows (DDoS, PortScan, Bot)
    are entirely off-diagonal.}
    \label{fig:cm}
\end{figure}

Per-class F1 on the day split is exactly 0.0 for DDoS, PortScan, and Bot
across all four models. The 0.44 macro F1 reported above comes entirely
from the Benign class and the few attack classes that overlap between
Mon--Wed training and Thu--Fri evaluation.

\section{Feature Importance and Per-Instance Interpretability}

The structural failure of multi-class classification on the day split
does not invalidate everything that has been learned. The features that
drive the in-distribution model tell us which flow statistics carry the
most discriminative power and therefore which features an attacker would
need to manipulate to evade detection. The XGBoost gain-based importance
ranks Backward Packet Length Std, Idle Mean, Forward Active Data Packets,
and Backward Packet Length Mean as the top four features. Server-side
response statistics carry most of the discriminative power. A
random-forest permutation-importance check produces a similar ranking,
which gives us confidence that the signal is real rather than a quirk
of the gain calculation.

\begin{figure}[H]
    \centering
    \includegraphics[width=0.8\textwidth]{feature_importance_xgboost.png}
    \caption{XGBoost gain-based feature importance (top features).}
    \label{fig:fi-xgb}
\end{figure}

For per-instance interpretability we use SHAP. The SHAP summary plot for
the binary random-forest detector shows that small Flow IAT Min values
push predictions strongly toward the attack class --- exactly the pattern
an attacker would observe and exploit, as confirmed in
Chapter~\ref{ch:adversarial}.

\begin{figure}[H]
    \centering
    \includegraphics[width=0.7\textwidth]{shap_summary_rf_binary.png}
    \caption{SHAP summary plot for the binary RF detector. Each dot is one
    flow; colour is feature value; position is the SHAP value.}
    \label{fig:shap-summary}
\end{figure}
